\chapter{The widget catalogue}
\label{ch:widgets}

The chapters so far described the engine. This chapter is a tour of
the cabin: the 69 widgets \maya{} ships under
\file{maya/include/maya/widget/}. We will not document every one ---
that is what \file{docs/13-widgets.md} and \file{docs/11-api-reference.md}
are for. Instead, we group them, explain the common interface, and
build a real composite UI to show how they cooperate.

\section{The widget contract}

A \maya{} widget is, by convention, a struct with:

\begin{itemize}[leftmargin=*]
  \item A \code{Config} struct of inputs (state to render).
  \item A free function \code{build(Config) -> Element} or a method
    \code{operator Element() const}.
  \item Optionally, a free \code{handle(Event\&, Config\&) ->
    std::optional<Msg>} for interactive widgets.
  \item Optionally, sub-widget-specific helper functions.
\end{itemize}

The widget never owns model state. It transforms a \code{Config}
into an \code{Element}. The host owns the state, threads it through
the widget on each render, and feeds events back. This is the same
discipline as \code{Program::view} but applied per-widget.

\section{Categories}

The widgets fall into five rough groups.

\subsection{Data display}

\begin{itemize}[leftmargin=*]
  \item \code{LineChart} --- a sparkline or full chart over time.
  \item \code{BarChart} --- horizontal or vertical bars.
  \item \code{Gauge} --- a circular percentage indicator.
  \item \code{Sparkline} --- a one-line trend indicator.
  \item \code{Heatmap} --- 2-D coloured grid for matrix data.
  \item \code{FlameChart} --- profiling-style stacked bars.
  \item \code{Waterfall} --- streaming spectrum or timing view.
  \item \code{TokenStream} --- the LLM ``tokens arriving live'' view.
  \item \code{ContextWindow} --- token budget bar.
  \item \code{GitGraph} --- commit tree.
\end{itemize}

These take ranges (\code{std::span<double>}, \code{std::vector<Sample>})
and produce a rectangular widget. Many use the canvas API
(Chapter~\ref{ch:canvas}) directly for sub-cell graphics (braille
characters, half-blocks).

\subsection{Input}

\begin{itemize}[leftmargin=*]
  \item \code{Input} --- single-line text field.
  \item \code{Textarea} --- multi-line editor with word wrap.
  \item \code{Select} --- single-choice dropdown.
  \item \code{Slider} --- numeric range picker.
  \item \code{Checkbox} --- boolean toggle.
  \item \code{Radio} --- mutually-exclusive choices.
  \item \code{Button} --- clickable rectangle.
  \item \code{Menu} --- vertical list of choices.
  \item \code{CommandPalette} --- a Ctrl-K style fuzzy launcher.
\end{itemize}

These widgets follow the \emph{controlled} pattern: their state
lives outside, in your model, and they receive it via \code{Config}.
On a relevant key, they return a \code{Msg} (or, in the quick-tool
API, mutate a signal).

\subsection{Layout helpers}

\begin{itemize}[leftmargin=*]
  \item \code{Table} --- rows and columns with sort indicators.
  \item \code{Tabs} --- horizontal tab strip + content area.
  \item \code{Tree} --- collapsible hierarchical view.
  \item \code{Scrollable} --- viewport over a too-tall child.
  \item \code{Modal} --- a floating dialog with backdrop.
  \item \code{Popup} --- a small contextual overlay.
  \item \code{Toast} --- a transient banner.
  \item \code{Disclosure} --- expand/collapse panel.
  \item \code{Divider} --- horizontal or vertical line.
  \item \code{Breadcrumb} --- path navigation crumbs.
\end{itemize}

\subsection{Display}

\begin{itemize}[leftmargin=*]
  \item \code{Markdown} --- a real (small) markdown renderer.
  \item \code{InlineDiff} --- single-line +/- diff view.
  \item \code{DiffView} --- multi-line diff with syntax highlight.
  \item \code{Badge} --- coloured pill (status, tag).
  \item \code{Spinner} --- animated activity indicator.
  \item \code{ProgressBar} --- discrete or continuous.
  \item \code{Calendar} --- month view.
  \item \code{Canvas} --- imperative drawing surface.
  \item \code{Image} --- iTerm2 / Kitty graphics protocol.
  \item \code{LogViewer} --- scrolling log tail.
\end{itemize}

\subsection{Agent UI (the \agentty{} family)}

Several widgets are designed for agentic chat UIs:

\begin{itemize}[leftmargin=*]
  \item \code{ToolCall}, \code{BashTool}, \code{ReadTool},
    \code{EditTool}, \code{WriteTool}, \code{FetchTool} --- a
    purpose-built card per tool kind.
  \item \code{Message}, \code{ThinkingBlock}, \code{StreamingCursor}.
  \item \code{ActivityBar}, \code{PermissionPrompt},
    \code{ModelBadge}, \code{CostTracker}.
  \item \code{GitStatus}, \code{FileChanges},
    \code{SystemBanner}, \code{ErrorBlock},
    \code{TurnDivider}, \code{ConversationView}, \code{PlanView},
    \code{ApiUsage}.
\end{itemize}

These are not part of any TUI orthodoxy --- they were built because
\agentty{} needed them. Their existence shows what the engine can
sustain. Their source is a worked tour of using the layout engine,
signals, and canvas together.

\section{Anatomy of a widget: \code{Input}}

Let us look at one widget in enough depth that you can read the
others. \code{Input} (file \file{widget/input.hpp}) is a single-line
text field.

\subsection{Configuration}

\begin{lstlisting}
struct InputConfig {
    std::string value;                  // current text
    int         cursor_pos = 0;          // character index
    std::string placeholder;             // dimmed text when empty
    int         max_length = 0;          // 0 = no limit
    bool        password = false;        // show as bullets
    Style       text_style{};            // optional override
    int         width = -1;              // -1 = grow
};
\end{lstlisting}

State out, view in.

\subsection{Building the element}

\begin{lstlisting}
Element input(const InputConfig& cfg);

inline Element input(const InputConfig& cfg) {
    auto display = cfg.value.empty()
        ? text(cfg.placeholder) | Dim
        : text(cfg.password ? std::string(cfg.value.size(), '*') : cfg.value);

    // The cursor block: a highlighted cell at cursor_pos.
    // (Real implementation handles word boundaries, scroll, etc.)
    return h(display);
}
\end{lstlisting}

The real implementation is fancier --- it scrolls horizontally when
the value exceeds the box width, it tracks the cursor as a separate
overlay, it respects bracketed paste. But the shape is: take
\code{Config}, return \code{Element}.

\subsection{Handling input}

\begin{lstlisting}
struct InputEvent {
    enum class Kind { Edited, Submitted, Cancelled };
    Kind        kind;
    std::string new_value;
    int         new_cursor;
};

std::optional<InputEvent> handle_input(const Event& ev, InputConfig& cfg);
\end{lstlisting}

The host calls \code{handle\_input(ev, my\_cfg)}. If the event is
relevant, the widget mutates \code{cfg} and returns an
\code{InputEvent} the host can translate into its own \code{Msg}.
If irrelevant, returns \code{nullopt}.

This pattern --- \code{handle(event, config) -> optional<event\_out>}
--- repeats across every interactive widget.

\section{Composing widgets}

A real screen is many widgets:

\begin{lstlisting}
auto chat_screen(const Model& m) -> Element {
    return v(
        h(
            // sidebar
            v(
                t<"Threads"> | Bold,
                map(m.threads, [](const auto& t) {
                    return text(t.title);
                })
            ) | width(20) | border_<Single>,

            // main conversation
            v(
                map(m.messages, [](const auto& msg) {
                    return message(MessageConfig{
                        .role  = msg.role,
                        .body  = msg.body,
                        .time  = msg.time
                    });
                })
            ) | grow_<1> | scrollable
        ) | grow_<1>,

        // composer at the bottom
        input(InputConfig{
            .value       = m.draft,
            .cursor_pos  = m.draft_cursor,
            .placeholder = "Message...",
        }) | border_<Round> | pad<0, 1>
    );
}
\end{lstlisting}

Each widget is a small function; the screen is composition. The
host's \code{update} routes events to whichever widget should
receive them, mutates the relevant config, returns the new model.

\section{Where the polish lives}

If you skim widget source files, you will notice they are mostly
\emph{boring}. Most of the implementation is layout, theming, and
small DSL chains. The interesting bits are concentrated in a few
widgets that exercise the canvas API (charts, sparklines,
heatmaps), the markdown renderer (a full incremental parser), and
the \code{Tool*} family (which model rich agent UIs).

The lesson: \maya{} is engineered so that adding a widget is a
short job. Picking a widget's name, defining a \code{Config},
writing a \code{build} function. The library you call once you have
read this book.

\begin{recap}
\maya{}'s widgets are small structs: a \code{Config} of inputs, a
\code{build} function returning \code{Element}, an optional
\code{handle} for events. They never own model state. The
catalogue spans data display, input, layout, display, and an
agent-UI family inherited from \agentty{}. The same pattern --- a
function from value to element --- scales from \code{Badge} to
\code{ConversationView}.
\end{recap}

\section*{Workshop}

\begin{enumerate}[leftmargin=*]
  \item Pick a widget you've never used. Read its header. Build a
    standalone example that drives it.
  \item Implement your own widget: a \code{Stopwatch} that takes a
    duration and renders \code{MM:SS.mmm}. Use the canvas API for
    a thick-stroked clock face.
  \item Read \file{maya/include/maya/widget/markdown.hpp}.
    Understand how it streams markdown incrementally without
    re-parsing the whole document on every keystroke.
\end{enumerate}
